% from ega2-1-fr.tex
\oldpage[II]{128}

Si on pose $h_\lambda=N(h'_\lambda)$, on aura donc les relations analogues
\[
N(\omega^{(2)}_{\lambda\mu})h_\mu
=h_\lambda N(\omega^{(1)}_{\lambda\mu})
\]
et par suite les $h_\lambda$ définissent un homomorphisme
$N(\mathcal{L}'_1)\to N(\mathcal{L}'_2)$ que nous désignerons par
$N_{\mathcal{B}/\mathcal{A}}(h')$ ou $N(h')$. Dans l'hypothèse (I), pour que
$h'$ soit un \emph{isomorphisme}, il faut et il suffit (puisqu'il s'agit d'une
question locale) que $N_{\mathcal{B}/\mathcal{A}}(h')$ le soit. Dans
l'hypothèse (II), cette condition est encore nécessaire~; elle est suffisante
si l'hypothèse (II \emph{bis}) est vérifiée et si
$\mathcal{B}\to\mathcal{B}\otimes_{\mathcal{A}}\mathcal{R}(X)$ est injectif.

Prenons en particulier $\mathcal{L}'_1=\mathcal{B}$~; les homomorphismes
$\mathcal{B}\to\mathcal{L}'$ s'identifient alors
(\hyperref[0.5.1.1-fr]{0, 5.1.1}) aux sections de $\mathcal{L}'$ au-dessus de
$X$, d'où une application canonique
\[
N_{\mathcal{B}/\mathcal{A}}:\Gamma(X,\mathcal{L}')
\to\Gamma(X,N_{\mathcal{B}/\mathcal{A}}(\mathcal{L}')).
\]
Il résulte encore de (\hyperref[II.6.5.1.1-fr]{6.5.1.1}) que si
$f'_1\in\Gamma(X,\mathcal{L}'_1)$, $f'_2\in\Gamma(X,\mathcal{L}'_2)$, on a
\[
\label{II.6.5.3.1-fr}
N(f'_1\otimes f'_2)=N(f'_1)\otimes N(f'_2).
\tag{6.5.3.1}
\]
Pour tout $\mathcal{A}$-Module inversible $\mathcal{L}$ et toute section
$f\in\Gamma(X,\mathcal{L})$, on a, compte tenu de
(\hyperref[II.6.5.2.4-fr]{6.5.2.4})
\[
\label{II.6.5.3.2-fr}
N_{\mathcal{B}/\mathcal{A}}(f\otimes1_{\mathcal{B}})=f^{\otimes n}
\tag{6.5.3.2}
\]
lorsque $\mathcal{B}$ est de rang constant $n$ dans l'hypothèse (I) (resp.
lorsque $\mathcal{B}\otimes_{\mathcal{A}}\mathcal{R}(X)$ est de rang constant
$n$ dans l'hypothèse (II)). Enfin, pour que l'homomorphisme
$\mathcal{B}\to\mathcal{L}'$ correspondant à une section $f'$ de
$\mathcal{L}'$ au-dessus de $X$ soit un isomorphisme, il faut et il suffit que
$f'_x$ soit une base de $\mathcal{L}'_x$ pour tout $x\in X$~; dans l'hypothèse
(I), cette condition équivaut donc à dire que $(N(f'))_x$ est une base de
$(N(\mathcal{L}'))_x$ pour tout $x$~; dans l'hypothèse (II), cette condition
est encore nécessaire, et elle est suffisante lorsque $\mathcal{B}$ vérifie
(II \emph{bis}) et que
$\mathcal{B}\to\mathcal{B}\otimes_{\mathcal{A}}\mathcal{R}(X)$ est injectif.
\end{env}

\begin{env}[6.5.4]
\label{II.6.5.4-fr}
Soient $(X,\mathcal{A})$, $(X',\mathcal{A}')$ deux espaces annelés,
$f:X'\to X$ un morphisme, $\mathcal{B}$ une $\mathcal{A}$-Algèbre,
$\mathcal{B}'=f^*(\mathcal{B})$ la $\mathcal{A}'$-Algèbre image réciproque. On
suppose vérifiée l'une des hypothèses suivantes~:
\begin{enumerate}
\item[$1^\circ$] $\mathcal{B}$ vérifie l'hypothèse (I) de
(\hyperref[II.6.5.1-fr]{6.5.1}).
\item[$2^\circ$] $(X,\mathcal{A})$ et $\mathcal{B}$ vérifient l'hypothèse
(II) de (\hyperref[II.6.5.1-fr]{6.5.1}), $(X',\mathcal{A}')$ est un préschéma
localement noethérien réduit, et si l'on désigne par $X_\alpha$ et
$X'_\beta$ les sous-préschémas fermés réduits respectifs de $X$ et $X'$ ayant
pour espaces sous-jacents les composantes irréductibles de ces espaces, la
restriction de $f$ à chaque $X'_\beta$ est un morphisme \emph{dominant} de
$X'_\beta$ dans un $X_\alpha$.
\end{enumerate}
Sous ces conditions, $\mathcal{B}'$ vérifie respectivement l'hypothèse (I) ou
l'hypothèse (II) de (\hyperref[II.6.5.1-fr]{6.5.1})~; la première assertion
est immédiate. Pour établir la seconde, il suffit de voir que pour tout
$x'\in X'$, les rangs des
$\mathcal{B}'\otimes_{\mathcal{O}_{X'}}\mathcal{R}(X'_\beta)$ pour les $\beta$
tels que $x'\in X'_\beta$ sont \emph{les mêmes}. Or, si la restriction de $f$
à $X'_\beta$ est un morphisme dominant dans $X_\alpha$, le rang de
$\mathcal{B}'\otimes_{\mathcal{O}_{X'}}\mathcal{R}(X'_\beta)$ est égal à celui
de $\mathcal{B}\otimes_{\mathcal{O}_X}\mathcal{R}(X_\alpha)$ (on le voit
aussitôt en se ramenant au cas affine comme dans
(\hyperref[II.6.4.9-fr]{6.4.9})), d'où notre assertion en vertu de l'hypothèse
(II) et de (\hyperref[II.6.5.1-fr]{6.5.1}).

